\documentclass[11pt,a4paper]{article}
\usepackage[utf8]{inputenc}
\usepackage[margin=0.8in]{geometry}
\usepackage{titlesec}
\usepackage{enumitem}
\usepackage{hyperref}

\hypersetup{
    colorlinks=true,
    linkcolor=black,
    filecolor=black,      
    urlcolor=black,
}

\titleformat{\section}{\Large\bfseries}{}{0em}{}[\titlerule]
\titlespacing{\section}{0pt}{1ex}{0.8ex}
\setlength{\parindent}{0pt}
\setlist[itemize]{leftmargin=*}

\begin{document}
\pagestyle{empty}

{\centering
\textbf{\huge Marcio Gabriel Schinor Mazega} \\ \vspace{0.2cm}
\textit{\large Desenvolvedor Backend e Full Cycle | Dados} \\ \vspace{0.1cm}
Tupã, SP -- Brasil | (14) 99812-6535 | \href{mailto:marciomazegaa@gmail.com}{marciomazegaa@gmail.com} \\
\href{https://linkedin.com/in/marcio-mazega}{linkedin.com/in/marcio-mazega} | \href{https://github.com/marciomazega}{github.com/marciomazega} \\
}

\vspace{0.3cm}
\section*{RESUMO PROFISSIONAL}
Desenvolvedor Full Cycle e Engenheiro de Backend com ampla visão do fluxo corporativo, da extração arquitetural de dados à interface com o usuário final. Experiência robusta na construção de APIs escaláveis (FastAPI, Spring Boot) e modelagem profunda de banco de dados (PostgreSQL), aliada ao domínio de frameworks Frontend (React, Angular). Capacidade comprovada de idealizar, estruturar e realizar o deploy de aplicações SaaS, soluções analíticas e rotinas de automação, garantindo alta performance e segurança desde a raiz tecnológica.

\vspace{0.1cm}
\section*{HABILIDADES TÉCNICAS}
\begin{itemize}[nosep]
    \item \textbf{Backend e APIs:} Python (FastAPI), Java (Spring Boot), Node.js, TypeScript, RESTful APIs.
    \item \textbf{Frontend e Web:} Angular, React, JavaScript, JWT Authentication, Mapas Interativos (Leaflet).
    \item \textbf{Bancos de Dados e Deploy:} PostgreSQL, MySQL, Docker, Linux (WSL).
    \item \textbf{Automação e Dados:} Scripts em Python, Integração de Sistemas, Web Scraping, Git / GitHub.
\end{itemize}

\vspace{0.1cm}
\section*{EXPERIÊNCIA PROFISSIONAL}
\textbf{Desenvolvedor Backend e Engenheiro de Dados} \hfill \textit{Fev 2024 -- Jun 2026} \\
\textit{Gattaz Health -- Remoto} \\
\begin{itemize}[nosep]
    \item Criação de APIs RESTful usando FastAPI para unificar ecossistemas analíticos e distribuir dados em larga escala.
    \item Modelagem física e lógica em banco de dados relacional (PostgreSQL), garantindo a persistência segura e performática dos sistemas backend e ferramentas de BI da companhia.
    \item Integração de fluxos de backend com sistemas de automação nativos em Python (Web Scraping / rotinas ETL).
\end{itemize}

\vspace{0.2cm}
\textbf{Pesquisador de Iniciação Científica (FAPESP/CNPq)} \hfill \textit{Fev 2025 -- Set 2025} \\
\textit{Smart B100 e OBA} \\
\begin{itemize}[nosep]
    \item Resolução de problemas de modelagem e mapeamento de sistemas interligados através de rigor científico.
\end{itemize}

\vspace{0.1cm}
\section*{PROJETOS DE DESTAQUE}
\textbf{Buscador Inteligente de Artigos}
\begin{itemize}[nosep]
    \item \textbf{Stack:} Python, FastAPI, Web Scraping.
    \item \textbf{Descrição:} Arquitetura de uma ferramenta automatizada para busca e extração de dados em fontes de publicações acadêmicas, conectando motores de scraping com uma API eficiente de resposta rápida.
\end{itemize}
\vspace{0.1cm}
\textbf{Ouvidoria Digital -- Synapse}
\begin{itemize}[nosep]
    \item \textbf{Stack:} Angular, Node.js, TypeScript, MySQL.
    \item \textbf{Descrição:} Solução Full Stack desenvolvida para Zeladoria Municipal que substitui processos em papel por mapas dinâmicos e sistema backend de autenticação (JWT) para gestão pública de chamados.
\end{itemize}
\vspace{0.1cm}
\textbf{Central Agro -- Plataforma SaaS}
\begin{itemize}[nosep]
    \item \textbf{Stack:} React, Python, GEE.
    \item \textbf{Descrição:} Pipeline backend de processamento geoespacial em Python entregue dinamicamente em interfaces frontend React modernas.
\end{itemize}

\vspace{0.1cm}
\section*{FORMAÇÃO ACADÊMICA E PREMIAÇÕES}
\begin{itemize}[nosep]
    \item \textbf{Graduação:} Tecnologia em Sistemas Inteligentes (IA e Dados) -- Fatec Shunji Nishimura (Dez 2027)
    \item \textbf{Premiação:} 3º Lugar -- DSIN Coder Challenge (Desafio de programação do Estado de SP, 2025)
    \item \textbf{Certificações:} Introdução ao Machine Learning (USP/ESALQ); Banco de Dados Relacional (EncData).
    \item \textbf{Publicações:} Coautor de 2 capítulos de livros acadêmicos sobre Agronegócio e Sociologia.
\end{itemize}

\end{document}